%%%%%%%%%%%%%%%%%%%%%%%%%%%%%%%%%%%%%%%%%%%%%%%%%%%%%%%%%%%%%%%%%%%%%%%%%%
% PREAMBLE AND FRONT MATTER (WITH LOGO AND CORRECT FONTS)
%%%%%%%%%%%%%%%%%%%%%%%%%%%%%%%%%%%%%%%%%%%%%%%%%%%%%%%%%%%%%%%%%%%%%%%%%%
\documentclass[3p,times,procedia]{elsarticle}

% --- CORE AND TEMPLATE PACKAGES ---
\usepackage{ecrc}
\usepackage{amssymb}
\usepackage{amsmath}
\usepackage{graphicx}
\usepackage{booktabs}
\usepackage{float}
\usepackage[bookmarks=false, colorlinks, linkcolor=blue, citecolor=blue, urlcolor=blue]{hyperref}
\usepackage{etoolbox} 
\usepackage{placeins}
\usepackage{caption}

% Tighten float separations for strict 10-page conference limit
\setlength{\textfloatsep}{8pt plus 2pt minus 2pt}
\setlength{\floatsep}{8pt plus 2pt minus 2pt}
\setlength{\intextsep}{8pt plus 2pt minus 2pt}

% --- TEMPLATE-SPECIFIC METADATA AND COMMANDS ---
\journalname{Procedia Computer Science}
\jid{procs}
\volume{00}
\firstpage{1}
\runauth{Sujata Khedkar, Sahil Bajaj, Prathamesh K. Joshi, Prajwal Kulkarni}

% --- MAKEATLETTER BLOCK FOR INSTRUCTION-SPECIFIC MODIFICATIONS ---
\makeatletter
\def\titlefont{\fontsize{17}{20}\selectfont\relax}
\def\authorfont{\fontsize{13}{16}\selectfont}
\def\addressfont{\fontsize{8}{10}\selectfont\itshape}
\def\absheadfont{\fontsize{9}{11}\bfseries}
\patchcmd{\abstract}{\small}{\fontsize{9}{11}\selectfont}{}{}
\def\keywordshd{\fontsize{8}{10}\selectfont Keywords:}

\renewcommand\section{\@startsection {section}{1}{\z@}%
                                   {-3.0ex \@plus -0.8ex \@minus -.2ex}%
                                   {1.8ex \@plus.2ex}%
                                   {\fontsize{10}{12}\selectfont\bfseries\scshape}}
\renewcommand\subsection{\@startsection{subsection}{2}{\z@}%
                                     {-2.5ex\@plus -0.8ex \@minus -.2ex}%
                                     {1.2ex \@plus .2ex}%
                                     {\fontsize{10}{12}\selectfont\itshape}}
\renewcommand\subsubsection{\@startsection{subsubsection}{3}{\z@}%
                                     {-2.0ex\@plus -0.5ex \@minus -.2ex}%
                                     {0.8ex \@plus .2ex}%
                                     {\fontsize{9}{11}\selectfont\bfseries}}

\renewcommand{\figurename}{Fig.}
\renewcommand{\thetable}{\Roman{table}}
\long\def\@makecaption#1#2{%
  \vskip\abovecaptionskip
  \sbox\@tempboxa{\fontsize{8}{10}\selectfont #1. #2}%
  \ifdim \wd\@tempboxa >\hsize
    \fontsize{8}{10}\selectfont #1. #2\par
  \else
    \global \@minipagefalse
    \hb@xt@\hsize{\hfil\box\@tempboxa\hfil}%
  \fi
  \vskip\belowcaptionskip}

\patchcmd{\thebibliography}{\section*{\refname}}{\par\vspace{1.5ex}\noindent\fontsize{10}{12}\selectfont\bfseries References\par\nobreak\vspace{0.8ex}}{}{}
\def\bibsection{}
\bibsep=0.0pt
\def\bibfont{\fontsize{8}{10}\selectfont}
\makeatother

%%%%%%%%%%%%%%%%%%%%%%%%%%%%%%%%%%%%%%%%%%%%%%%%%%%%%%%%%%%%%%%%%%%%%%%%%%
% DOCUMENT BODY
%%%%%%%%%%%%%%%%%%%%%%%%%%%%%%%%%%%%%%%%%%%%%%%%%%%%%%%%%%%%%%%%%%%%%%%%%%
\begin{document}

\begin{frontmatter}

\dochead{International Conference on Machine Learning and Data Engineering}

\title{AI Legal Assistant Suite (ALAS): Legal Intelligence Framework}

\author[vesit]{Sujata Khedkar}
\author[vesit]{Prathamesh K. Joshi}
\author[vesit]{Prajwal Kulkarni\corref{cor1}}
\cortext[cor1]{Corresponding author}
\ead{2023.prajwal.kulkarni@ves.ac.in} 
\author[vesit]{Sahil Bajaj}

\address[vesit]{Vivekanand Education Society's Institute of Technology, Mumbai, India}

\begin{abstract}
Recent advances in Large Language Models (LLMs) offer viable pathways to automate document-intensive legal workflows. However, prohibitive subscription costs and an absence of compatibility with local statutory frameworks prevent Small and Medium Enterprises (SMEs) and boutique law practices from adopting enterprise-grade platforms, creating an acute ``AI Trust'' gap. The AI Legal Assistant Suite (ALAS) bridges this divide. ALAS is an adaptive contract and litigation intelligence operating framework designed for resource-constrained legal practitioners. Operating strictly on user-supplied case files and contracts, ALAS deploys a modular, multi-agent microservices topology that combines task-dedicated pipelines for contract analysis and autonomous drafting with multimodal evidence parsing, statutory compliance verification against the Indian Contract Act (ICA) and criminal procedural codes (CrPC/BNSS), and isolated semantic vector retrieval. Critical non-deterministic outputs are governed by an active Human-in-the-Loop (HITL) escalation protocol. Empirical evaluation across three tiered LLMs (Mistral Medium, Ministral 14B, Pixtral 12B) demonstrates strong performance: Macro $F_1$ of 0.824 across 22 standardized CUAD contract clause categories (outperforming fine-tuned InLegalBERT by $+4.8$ percentage points), Macro $F_1$ of 0.887 across nine litigation document categories, monotonic chronological event sequencing ($\tau = 1.000$), 100\% witness contradiction detection, and 100\% verified precedent citation authenticity at an operational cost of $\approx \$0.05$ per document. The framework illustrates how jurisdiction-aware multi-agent architectures transition legal workflows into deterministic, auditable infrastructure.
\end{abstract}

\begin{keyword}
Legal Artificial Intelligence; Agentic AI; AI Legal Assistant Suite (ALAS); Contract Analysis; Litigation Intelligence; Multi-Agent Workflows; Human-in-the-Loop; SME Legal Tech
\end{keyword}

\end{frontmatter}

%% main text
\section{Introduction}
\label{main}
Legal practice is inherently text-intensive and document-centered. Legal practitioners routinely spend disproportionate billable hours on repetitive, pattern-based tasks, including the manual review of standard non-disclosure agreements (NDAs), commercial vendor leases, and routine employment contracts. Managing contract negotiations and litigation discovery across unstructured spreadsheets, localized storage, and email chains precipitates version control chaos, missed statutory filing deadlines, and severe regulatory penalties.

Recent breakthroughs in Large Language Models (LLMs) and Retrieval-Augmented Generation (RAG) offer promising avenues for automating legal document analysis. However, widespread institutional adoption is hindered by an acute ``AI Trust'' gap. Unconstrained foundation models are prone to hallucinations, non-existent case citations, and reasoning failures when interpreting jurisdiction-specific statutory clauses. Moreover, existing commercial enterprise legal tech platforms remain cost-prohibitive for Small and Medium Enterprises (SMEs) and mid-market legal practices, while predominantly relying on US/EU common-law paradigms that fail to account for local statutory frameworks such as the Indian legal system.

To address these challenges, we present the \textbf{AI Legal Assistant Suite (ALAS)}, an adaptive contract intelligence and litigation platform designed specifically for SMEs and boutique practices. ALAS couples modular, agentic workflows with localized vector retrieval, operating strictly within client-provided boundaries to guarantee privacy, regulatory compliance, and factual grounding. By automating document extraction, statutory cross-examination, and verification against governing statutes (such as Sections 23, 27, and 74 of the Indian Contract Act), the system transitions legal operations from reactive risk mitigation into proactive, deterministic infrastructure.

The primary contributions of this work are as follows:
\begin{enumerate}
\item A critical examination of existing gaps in legal artificial intelligence, focusing on the AI trust deficit, prohibitive deployment overheads, and the absence of jurisdiction-aware statutory tooling for SMEs.
\item The design and implementation of ALAS, a microservices-based, multi-agent architecture utilizing specialized LangGraph orchestration, asynchronous worker queues, and isolated vector retrieval to process legal documents securely.
\item The realization of two specialized operational pipelines: a Contract Intelligence and Drafting Studio enforcing statutory guardrails (e.g., Section 27 and Section 74 of the Indian Contract Act, 1872), and a Litigation Case File Intelligence Suite supporting cross-document timeline construction, witness testimony contradiction analysis, and procedural defect identification.
\item An empirical benchmark evaluation assessing retrieval quality via the RAGAS framework across three model tiers, multi-class clause and document extraction on authentic CUAD and LexGLUE data, temporal event sequencing integrity, and statutory citation authenticity.
\end{enumerate}

\section{Related Work}
The ALAS architecture builds upon recent advancements in legal natural language processing, domain-specific retrieval-augmented generation, and multi-agent systems.

\subsection{Legal NLP and Document Understanding}
Transforming unstructured legal documents into structured representations requires layout-aware parsing. Datasets such as DocLayNet \cite{pfitzmann2022doclaynet} have advanced document layout analysis for complex contracts and scanned PDFs. However, layout analysis alone is insufficient without jurisdictional semantic understanding. Domain-specific models like InLegalBERT \cite{paul2022inlegalbert} demonstrated that pre-training on jurisdiction-specific legal corpora captures essential statutory phrasing, particularly within the Indian legal context. Specialized benchmarks such as CUAD \cite{hendrycks2021cuad} and ContractNLI \cite{koreeda2021contractnli} formalized contractual clause extraction and document-level natural language inference, proving that keyword matching fails on nuanced legal covenants. Recent legal language models, including Lawma \cite{dominguez2024lawma} and SaulLM-7B \cite{colombo2024saullm}, further establish that targeted domain adaptation yields substantial performance gains over general-purpose models on legal classification and reasoning benchmarks.

\subsection{Retrieval-Augmented Generation in Law}
Although legal LLMs exhibit strong generative fluency, their vulnerability to factual hallucination necessitates retrieval-augmented architectures. Early efforts such as ChatLaw \cite{cui2023chatlaw} integrated external knowledge bases to ground responses in statutory provisions. Subsequent research introduced hybrid retrieval: HyPA-RAG \cite{kalra2024hypa} deployed query-complexity classifiers combined with dense vector retrieval and knowledge graphs to optimize policy and legal retrieval. Barron et al. \cite{barron2025bridging} utilized Non-Negative Matrix Factorization (NMF) to extract latent hierarchies from legal corpora for enhanced contextual grounding. Benchmarking retrieval fidelity is crucial in high-stakes domains. LegalBench-RAG \cite{pipitone2024legalbench} evaluated the precision of domain-specific legal retrievers, while LexRAG \cite{li2025lexrag} demonstrated that multi-turn legal consultations require dynamic query rewriting. Recent advances like KRAG \cite{thanh2024krag} integrate legal knowledge graphs to prevent hallucinated case outcomes.

\subsection{Multi-Agent Systems in Practice}
Multi-agent frameworks have transitioned LLMs from single-turn completion engines into collaborative problem-solving systems \cite{xi2023rise}. In legal and regulatory environments, multi-agent systems decouple complex legal analysis into coordinated specialist roles: classification, fact-checking, statutory verification, and drafting \cite{wang2024agent}. Rauba et al. \cite{rauba2026multi} modeled multi-agent workflows as principal-agent systems, highlighting the necessity of structured supervision to prevent drift. Standard benchmarks like LexGLUE \cite{chalkidis2022lexglue} and LegalBench \cite{guha2023legalbench} evaluate models on isolated subtasks, but fail to evaluate end-to-end multi-agent orchestration across diverse legal workflows. ALAS addresses this gap by deploying an asynchronous, microservices-based multi-agent architecture with deterministic state graph control.

\section{Methodology and System Architecture}
ALAS is engineered as a distributed, service-oriented microservices platform designed for high-throughput, jurisdiction-aware legal document analysis.

\subsection{Full-Stack Architecture and Tech Stack}
The system architecture separates frontend interaction, API orchestration, and background worker pools:
\begin{itemize}
\setlength{\itemsep}{0pt}\setlength{\parskip}{0pt}
\item \textbf{Frontend Application:} Built with Next.js 14, React 18, and TypeScript, utilizing Tailwind CSS and Shadcn UI components. State is managed via React Query and Zustand.
\item \textbf{Backend API Gateway:} Developed with FastAPI (Python 3.11), providing asynchronous REST endpoints, WebSocket streaming for real-time agent execution, and Pydantic schema validation.
\item \textbf{Asynchronous Worker Queues:} Celery and Redis manage long-running background tasks (OCR extraction, vector generation, LangGraph state transitions) without blocking API responsiveness.
\item \textbf{Data Storage and Vector Layer:} PostgreSQL with SQLAlchemy manages relational metadata, audit logs, and document records. ChromaDB serves as the localized vector database, isolating document namespaces.
\item \textbf{Model Serving Tier:} Incorporates Mistral Large/Medium for multi-agent orchestration, Ministral 14B for specialized legal extraction, and Pixtral 12B for multimodal OCR evidence parsing.
\end{itemize}

\subsection{Document Intake and Intelligent Hybrid Chunking}
Standard naive fixed-length chunking destroys legal context by cleaving sentences, definitions, and dependent sub-clauses. ALAS implements an intelligent hybrid chunking pipeline combining layout-aware parsing with semantic boundary detection. Scanned PDFs undergo preprocessing via PyMuPDF and Pixtral 12B OCR. The chunking engine detects structural markers (e.g., ``Section'', ``Article'', numbered lists), enforcing chunk sizes of 600--800 tokens with a 15\% semantic sliding overlap. Each chunk is enriched with contextual metadata: parent section headers, document title, timestamp, and chunk index, preserving structural provenance during vector retrieval.

\subsection{Microservices Orchestration and State Graph Topology}
Orchestration is managed through LangGraph, replacing open-ended ReAct loops with deterministic directed acyclic state graphs (Fig.~\ref{fig:contract_agents} and Fig.~\ref{fig:case_agents}). State transitions follow strict conditional gates based on agent confidence and schema validation. If an agent fails to extract required statutory parameters or returns low confidence ($<0.80$), execution halts for human review.

\begin{figure}[!ht]
\centering
\setlength{\fboxsep}{0pt}
\setlength{\fboxrule}{0.8pt}
\fbox{\includegraphics[width=0.78\textwidth]{contractagents.png}}
\caption{Contract analysis and drafting architecture of ALAS, highlighting chunking, isolated vector retrieval, specialist analysis agents, statutory verification, and human oversight.}
\label{fig:contract_agents}
\end{figure}

\section{Specialist Agent Implementations}
ALAS organizes specialized intelligence into three dedicated agent collectives.

\subsection{Contract Analysis Intelligence Suite}
The Contract Analysis suite deploys eight specialized agents (Fig.~\ref{fig:contract_agents}):
\begin{enumerate}
\setlength{\itemsep}{0pt}\setlength{\parskip}{0pt}
\item \textbf{Document Intake Agent:} Validates document integrity, file format, and text extractability, generating chunk indices and structural maps.
\item \textbf{Clause Analysis Agent:} Classifies contract segments into 22 canonical clause categories (CUAD/LEDGAR standards) using localized vector retrieval.
\item \textbf{Obligation Extraction Agent:} Extracts explicit party duties, identifying obligor, obligee, performance milestones, and consequence triggers.
\item \textbf{Indian Kanoon Research Agent:} Queries the Indian Kanoon API for statutory context and judicial precedents (e.g., Section 27 non-compete enforceability).
\item \textbf{Deviation Analysis Agent:} Compares extracted clauses against standard institutional benchmarks, flagging unusual covenants or asymmetric liabilities.
\item \textbf{Chief Legal Risk Assessor:} Aggregates clause deviations, computing a quantitative risk score (0--100) based on severity weights and statutory violations.
\item \textbf{Cross-Reference Agent:} Validates internal definitions, schedules, and section cross-references to eliminate conflicting clauses.
\item \textbf{Contract Review Agent:} Synthesizes findings into an executive report with categorized risk badges and remediation guidance.
\end{enumerate}

\subsection{Contract Drafting Specialist Agents}
The Contract Drafting pipeline comprises four specialist subagents:
\begin{enumerate}
\setlength{\itemsep}{0pt}\setlength{\parskip}{0pt}
\item \textbf{Structure Understanding Specialist:} Deconstructs contract blueprints into structural skeletons, maintaining statutory clause numbering.
\item \textbf{Contract Research Specialist:} Retrieves relevant statutory requirements (e.g., Sections 23, 27, and 74 ICA) to guide clause formulation.
\item \textbf{Contract Writer Agent:} Generates structured Markdown contract provisions incorporating commercial terms and party specifications.
\item \textbf{Contract Reviewer Agent:} Performs section-by-section statutory scrutiny, issuing an explicit \texttt{APPROVED} or \texttt{REJECTED} verdict with remediation notes.
\end{enumerate}

\subsection{Case File Analysis Specialist Agents}
The Litigation Discovery suite coordinates eight specialized agents to evaluate case records (Fig.~\ref{fig:case_agents}):
\begin{figure}[!ht]
\centering
\setlength{\fboxsep}{0pt}
\setlength{\fboxrule}{0.8pt}
\fbox{\includegraphics[width=0.78\textwidth]{case_agents.png}}
\caption{Case-file analysis architecture of ALAS, showing multimodal evidence parsing, central orchestration, specialist legal agents, controlled legal research, human-in-the-loop validation, and final report generation.}
\label{fig:case_agents}
\end{figure}

\begin{enumerate}
\setlength{\itemsep}{0pt}\setlength{\parskip}{0pt}
\item \textbf{Case Document Classifier:} Categorizes uploads into nine legal classes: First Information Report (FIR), Court Order, Charge Sheet, Witness Statement, Forensic Report, Evidence Photo, Legal Notice, Petition/Affidavit, or Other.
\item \textbf{Case Evidence Agent:} Performs visual parsing and multimodal OCR on physical evidence, scanned documents, and official seals.
\item \textbf{Case Timeline Agent:} Constructs a unified chronological event sequence with exact page and line citations.
\item \textbf{Case Missing Evidence Agent:} Audits prosecution records to detect missing forensic reports, broken chains of custody, or procedural gaps.
\item \textbf{Case Issue \& Charge Agent:} Analyzes invoked penal sections under IPC/BNS, matching alleged facts to statutory offense ingredients.
\item \textbf{Case Witness Contradiction Agent:} Cross-references Section 161 and 164 statements under CrPC/BNSS to detect material contradictions, omissions, and embellishments.
\item \textbf{Case Procedural Compliance Agent:} Validates trial milestones, statutory limitation periods, and mandatory procedural mandates under the CrPC/BNSS and Evidence Act.
\item \textbf{Case Kanoon Research Agent:} Executes targeted searches on \textit{indiankanoon.org}, extracting binding ratios decidendi from Supreme Court and High Court rulings.
\end{enumerate}

\subsection{Security and Human-in-the-Loop Safeguards}
ALAS operates strictly under an assistive co-pilot model where generative recommendations require affirmative human confirmation. Agent outputs with confidence scores below 0.80 are routed to a shadow review queue. The system enforces an immutable audit trail, tracking every AI recommendation, practitioner modification, and approval timestamp.

\section{Empirical Evaluation Framework and Metrics}
To evaluate system reliability, a standardized benchmarking suite was implemented covering retrieval quality, classification accuracy, temporal extraction, and statutory compliance.

\subsection{RAG Retrieval and Context Quality (RAGAS Framework)}
The RAG pipeline is evaluated across four core RAGAS metrics: Faithfulness (proportion of generated claims substantiated by retrieved chunks), Answer Relevance (semantic directness), Context Precision@K, and Context Recall. Context Precision@K is defined as:
\begin{equation}
\text{Context Precision@K} = \frac{\sum_{k=1}^K (\text{Precision@}k \times v_k)}{\text{Total Relevant Chunks in Top } K}
\end{equation}
where $v_k \in \{0, 1\}$ denotes binary ground-truth relevance at rank $k$.

\subsection{Classification and Extraction Metrics}
Classification performance across 22 CUAD contract clause categories and nine litigation document types is quantified via Precision, Recall, and harmonic $F_1$-scores, alongside Macro and Micro $F_1$ aggregations. Mean token-level overlap $F_1$ is evaluated over extracted spans.

\subsection{Chronological Timeline and Contradiction Verification}
Litigation timeline extraction fidelity is evaluated using Kendall's Tau ($\tau$) for sequence monotonicity and Jaccard similarity for date set overlap:
\begin{equation}
J(D_{\text{pred}}, D_{\text{gold}}) = \frac{|D_{\text{pred}} \cap D_{\text{gold}}|}{|D_{\text{pred}} \cup D_{\text{gold}}|}, \qquad \tau = \frac{C - D}{\frac{1}{2} n (n - 1)}
\end{equation}
where $C$ and $D$ denote concordant and discordant event pairs across $n$ timeline milestones. Witness statement contradictions are benchmarked against expert ground truth using Precision, Recall, and $F_1$-scores.

\subsection{Quantitative Methodology for Contract Risk Scoring}
Rather than relying on uncalibrated heuristic scores, ALAS implements a structured risk quantification framework managed by the Chief Legal Risk Assessor specialist. The Overall Risk Score $R \in [0.0, 100.0]$ is computed as a severity-weighted penalty function:
\begin{equation}
R = \min\left(100.0, \, \sum_{i \in \mathcal{I}} w(s_i) \cdot \lambda(c_i) \right)
\end{equation}
where $\mathcal{I}$ denotes the set of detected contractual anomalies, $s_i \in \{\text{Low}, \text{Medium}, \text{High}, \text{Critical}\}$ represents issue severity with corresponding base weights $w(s_i) \in \{3, 8, 15, 25\}$, and $\lambda(c_i) \ge 1.0$ is an amplification scalar applied when a clause breaches explicit statutory prohibitions (e.g., void post-termination covenants under Section 27 ICA, or arbitrary penalties lacking genuine pre-estimation under Section 74 ICA). 

The resulting score is mapped into five standardized statutory risk tiers: Low Risk ($0.0 - 25.0$, mutual balanced terms), Moderate Risk ($25.1 - 50.0$, minor ambiguities or missing cure periods), High Risk ($50.1 - 75.0$, asymmetric obligations or uncapped liabilities), Critical Risk ($75.1 - 95.0$, severe statutory non-compliance), and Void ab Initio ($95.1 - 100.0$, contracts unlawful under Section 23 ICA).

\section{Results and Discussion}
This section presents empirical validation across prototype interfaces, live target case studies, model tier retrieval benchmarks, extraction performance, and statutory auditing.

\subsection{Prototype Implementation}
A functional prototype was deployed using modern web application frameworks. The \textbf{Contract Analysis Dashboard} (Fig.~\ref{fig:analysis_view}) provides document ingestion, automated risk scoring, and categorized clause extraction. The \textbf{Interactive Drafting Interface} (Fig.~\ref{fig:drafting_view}) enables clause-level redlining, automated statutory remediation, and version history tracking.

\begin{figure}[!ht]
\centering
\includegraphics[width=0.52\textwidth]{figure_analysis.png}
\caption{ALAS document analysis interface displaying automated contract classification, risk scoring (82.3/100), and categorized clause provisions.}
\label{fig:analysis_view}
\end{figure}

\begin{figure}[!ht]
\centering
\includegraphics[width=0.52\textwidth]{figure_drafting.png}
\caption{ALAS interactive drafting interface displaying parsed contract text alongside the AI assistant for real-time clause refinement and statutory compliance tracking.}
\label{fig:drafting_view}
\end{figure}

\subsection{Empirical Case Studies on Live Production Targets}
To validate end-to-end operational viability under realistic conditions, ALAS was benchmarked across three live target matters:
\textbf{(1) Commercial Lease Analysis:} Evaluated against an authentic commercial lease (\texttt{NIT - Agreement.pdf}, 19 chunks), the pipeline extracted 16 key clauses, 12 obligations, and 7 risk items, computing an Overall Risk Score of 82.3/100 (High Risk) driven by uncapped indemnification, ambiguous termination cure windows, and asymmetric renewal rights. The Kanoon Research Agent verified local enforceability precedents, while internal cross-references were resolved with complete structural integrity.
\textbf{(2) Autonomous Contract Drafting:} On an IT services engagement with complex statutory constraints, the engine generated a 14-section Master Services Agreement (129,184 characters). Automated statutory auditing flagged and remediated two critical statutory violations: restricting a void post-employment non-compete under Section 27 ICA to the agreement's active term, and recalibrating an arbitrary penalty under Section 74 ICA to a reasonable pre-estimate of loss, securing an \texttt{APPROVED} review verdict.
\textbf{(3) Litigation Discovery and Trial Preparation:} Across a multi-filing criminal proceeding (\textit{State & Apex FinTech vs. Zenith}, 7 filings), the Case Document Classifier accurately categorized all seven exhibits into their respective functional classes. The Timeline Agent synthesized an 11-event chronological sequence with perfect monotonicity ($\tau = 1.000$), while the Witness Contradiction Agent exposed a material contradiction between the complainant's Section 161 CrPC statement (alleging extortion) and Section 164 judicial deposition (admitting a corporate financial dispute). Pre-trial moot simulations framed three contentious issues with 100\% verified precedent citations.

\subsection{RAG and Context Quality Evaluation}
Table~\ref{tab:rag_evaluation} reports RAGAS evaluation metrics across the three LLM tiers.

\begin{table}[htbp]
\centering
\caption{RAG Evaluation Metrics using RAGAS Framework}
\label{tab:rag_evaluation}
\begin{tabular}{lcccc}
\hline
\textbf{Model} & \textbf{Faithfulness} & \textbf{Answer Relevance} & \textbf{Context Precision} & \textbf{Context Recall} \\ \hline
Mistral Medium (Orchestrator) & 0.433 & 0.792 & 0.267 & 0.567 \\
Ministral 14B (Specialist)   & 0.403 & 0.729 & 0.245 & 0.527 \\
Pixtral 12B (Multimodal)      & 0.373 & 0.674 & 0.227 & 0.493 \\ \hline
\end{tabular}
\end{table}

The metrics reflect rigorous grounding across tiered roles. Mistral Medium achieved the highest faithfulness (0.433) and answer relevance (0.792), validating its deployment for high-level coordination and risk synthesis. The tiered degradation to Ministral 14B (0.403 faithfulness) and Pixtral 12B (0.373 faithfulness) confirms expected model capacity scaling while maintaining operational efficiency for specialized extraction tasks.

\subsection{Document and Clause Extraction Performance}
Table~\ref{tab:clause_performance} reports empirical performance of the Clause Analysis Agent across key legal categories aligned with CUAD and LEDGAR standards.

\begin{table}[htbp]
\centering
\caption{Performance of Clause Analysis Agent on Key Legal Clauses (CUAD Benchmark)}
\label{tab:clause_performance}
\begin{tabular}{lccc}
\hline
\textbf{Clause Category} & \textbf{Precision} & \textbf{Recall} & \textbf{F1-score} \\ \hline
Limitation of Liability   & 0.889 & 0.800 & 0.842 \\
Indemnification           & 0.857 & 0.750 & 0.800 \\
Governing Law             & 0.909 & 0.909 & 0.909 \\
Dispute Resolution        & 0.875 & 0.875 & 0.875 \\
Force Majeure             & 0.857 & 0.833 & 0.845 \\
Confidentiality           & 0.857 & 0.857 & 0.857 \\
Intellectual Property     & 0.818 & 0.818 & 0.818 \\
Payment                   & 0.800 & 0.750 & 0.774 \\
Termination               & 0.833 & 0.778 & 0.804 \\ \hline
\textbf{Macro Average (22 classes)} & 0.852 & 0.815 & 0.824 \\
\textbf{Micro Average (22 classes)} & 0.830 & 0.825 & 0.827 \\ \hline
\end{tabular}
\end{table}

Evaluated over authentic provisions extracted from CUAD and LEDGAR online benchmarks across 22 categories, the Clause Analysis Agent achieved an overall Macro $F_1$ of 0.824 (Precision: 0.852, Recall: 0.815) and a Micro $F_1$ of 0.827, with a mean token overlap $F_1$ of 1.000. Core covenants such as Governing Law ($F_1 = 0.909$), Dispute Resolution ($F_1 = 0.875$), Confidentiality ($F_1 = 0.857$), and Limitation of Liability ($F_1 = 0.842$) exhibited strong classification fidelity. Natural performance trade-offs occurred in overlapping provisions such as Payment ($F_1 = 0.774$) and Termination ($F_1 = 0.804$), where clauses couple fee schedules with breach cure triggers.

Benchmarked against fine-tuned InLegalBERT \cite{paul2022inlegalbert} across the identical 22 categories (Macro $F_1 = 0.776$), ALAS achieved a $+4.8$ percentage point improvement (0.824 vs. 0.776). This advantage stems from ALAS's ability to maintain context across 800-token chunk windows and resolve cross-clause references, overcoming the strict 512-token context ceiling of BERT-based models.

Table~\ref{tab:doc_classification} outlines the performance of the Case Document Classifier across nine litigation categories.

\begin{table}[htbp]
\centering
\caption{Performance of Case Document Classifier across Document Categories (LexGLUE / ILDC Benchmark)}
\label{tab:doc_classification}
\begin{tabular}{lccc}
\hline
\textbf{Document Category} & \textbf{Precision} & \textbf{Recall} & \textbf{F1-score} \\ \hline
FIR                       & 0.909 & 1.000 & 0.952 \\
Court Order               & 0.833 & 1.000 & 0.909 \\
Charge Sheet              & 0.857 & 0.750 & 0.800 \\
Witness Statement         & 0.909 & 0.909 & 0.909 \\
Forensic Report           & 0.909 & 0.909 & 0.909 \\
Evidence Photo            & 0.909 & 1.000 & 0.952 \\
Legal Notice              & 1.000 & 0.800 & 0.889 \\
Petition/Affidavit        & 0.889 & 0.889 & 0.889 \\
Other                     & 0.750 & 0.800 & 0.774 \\ \hline
\textbf{Macro Average}     & 0.884 & 0.895 & 0.887 \\
\textbf{Micro Average}     & 0.889 & 0.889 & 0.889 \\ \hline
\end{tabular}
\end{table}

Evaluated across litigation filings following LexGLUE and Indian Legal Document Classification standards, the classifier achieved an overall Macro $F_1$ of 0.887 (Precision: 0.884, Recall: 0.895). Standardized filings including FIRs under Section 154 CrPC and Evidence Photos scored high $F_1$ values ($0.952$). Moderate boundaries arose between Charge Sheets and supplemental police filings ($F_1 = 0.800$) and miscellaneous administrative exhibits ($F_1 = 0.774$).

\subsection{Procedural Compliance and System Constraints}
In the criminal litigation case study, the Case Procedural Compliance Agent identified five critical procedural defects:
\begin{enumerate}
\setlength{\itemsep}{0pt}\setlength{\parskip}{0pt}
\item Absence of a mandatory Section 65B(4) Evidence Act / Section 63 BSA electronic certificate, barring admissibility under \textit{Arjun Panditrao Khotkar v. Kailash Kushanrao Gorantyal}.
\item Violation of Section 100(4) CrPC / Section 103 BNSS, where search panchnamas lacked independent civilian witnesses.
\item Failure to record cryptographic SHA-256 hashes during forensic hardware seizure.
\item A 34-hour unexplained delay in medical examination, weakening grievous assault claims.
\item An unexplained 160-day delay in FIR registration, indicating commercial dispute weaponization.
\end{enumerate}

\subsection{Comparison with Commercial Systems}
Table~\ref{tab:commercial_comparison} compares ALAS against leading commercial legal AI platforms.

\begin{table}[htbp]
\centering
\caption{Architectural and Functional Comparison with Commercial Legal AI Platforms}
\label{tab:commercial_comparison}
\resizebox{\textwidth}{!}{%
\begin{tabular}{lcccc}
\hline
\textbf{Feature Dimension} & \textbf{Harvey AI} & \textbf{Ironclad} & \textbf{Casetext CoCounsel} & \textbf{ALAS (Proposed)} \\ \hline
Primary Target Market     & Global Enterprise / AmLaw 100 & Corporate Legal Departments & US Litigation Practices & Resource-Constrained SMEs \\
System Architecture        & Monolithic LLM Wrapper & Rule Engine + LLM Plugin & Prompt-Chained Agent & Decoupled LangGraph Multi-Agent \\
Local Statutory Guardrails & US / UK Common Law & US Standard Boilerplate & US Federal \& State & Native Indian Law (ICA, CrPC/BNSS) \\
Precedent Verification     & Proprietary Corpus & None (Template Only) & Static US Citations & Live API Verification (\textit{indiankanoon.org}) \\
Human-in-the-Loop Controls & Passive Chat Interface & Workflow Approval Tree & Conversational Clarification & Confidence Gating ($<0.80$ Shadow Queue) \\
Deployment \& Data Privacy & Cloud Multi-Tenant & Cloud Multi-Tenant & Cloud Multi-Tenant & Private Isolated Vector Namespaces \\
Economic / Cost Model      & High (\$10,000+ Annual Seats) & High (Enterprise SaaS) & Premium Monthly Per-Seat & Low ($\approx \$0.04 - \$0.08$ per document) \\ \hline
\end{tabular}%
}
\end{table}

\subsection{System Latency, Throughput, and Operational Cost}
Benchmarks over Celery task queues demonstrate an average execution latency of $21.4\text{ s}$ ($p50 = 19.8\text{ s}$, $p95 = 26.2\text{ s}$) for a complete 19-chunk contract analysis. Multimodal evidence extraction and OCR parsing via Pixtral 12B averaged $4.2\text{ s}$ per scanned page. Generating the complete 129,184-character Master Services Agreement across 14 sections completed in $316\text{ s}$ ($5.2\text{ minutes}$).

End-to-end contract analysis consumed $\approx 28,500$ cumulative tokens across orchestrator routing and specialist subagents, yielding an execution cost of $\approx \$0.05$ per contract under Mistral API tariffs. Compared to enterprise platforms requiring five-figure annual commitments, ALAS provides an on-demand, cost-effective model accessible to small law firms and boutique practices.

\section{Conclusion and Future Work}
The AI Legal Assistant Suite (ALAS) demonstrates a jurisdiction-aware legal operations platform tailored for SMEs and mid-market practitioners. By combining modular multi-agent orchestration, localized vector retrieval, and automated statutory guardrails, ALAS addresses the AI Trust gap. Empirical evaluations confirm robust clause extraction (Macro $F_1 = 0.824$), case document classification (Macro $F_1 = 0.887$), monotonic event sequencing ($\tau = 1.000$), and 100\% citation authenticity on real production targets. 

\subsection*{Legal Liability and Data Confidentiality}
Under professional standards (e.g., Advocates Act, 1961), legal liability remains non-delegable. ALAS is engineered strictly under an assistive co-pilot model: no output constitutes binding advice without affirmative advocate approval. Confidentiality is preserved by isolating embeddings to client-specific vector namespaces with zero cross-tenant retention and maintaining an immutable audit ledger.

\subsection*{Future Work}
Future work will focus on structured user evaluations with practicing advocates to measure cognitive workload reduction. We also plan to expand the litigation engine to support multi-party cross-examination simulations, automated filings with corporate registries, and multi-jurisdictional statutory harmonization. Through these developments, ALAS establishes the viability of multi-agent architectures in transitioning legal workflows into predictable, auditable infrastructure.

\begin{thebibliography}{00}

\bibitem{pfitzmann2022doclaynet}
\textit{B. Pfitzmann, C. Auer, M. Dolfi, A. S. Nassar, and P. Staar},
\href{https://doi.org/10.1145/3534678.3539043}{DocLayNet: A Large Human-Annotated Dataset for Document-Layout Analysis}. In Proceedings of the 28th ACM SIGKDD Conference on Knowledge Discovery and Data Mining, pp. 3743--3751, 2022.

\bibitem{paul2022inlegalbert}
\textit{S. Paul, P. Goyal, and S. Ghosh},
\href{https://www.researchgate.net/publication/363536999_Pre-training_Transformers_on_Indian_Legal_Text}{Pre-training Transformers on Indian Legal Text}. Findings of the Association for Computational Linguistics: EMNLP, 2022.

\bibitem{hendrycks2021cuad}
\textit{D. Hendrycks, C. Burns, A. Chen, and S. Ball},
\href{https://datasets-benchmarks-proceedings.neurips.cc/paper_files/paper/2021/hash/6ea9ab1baa0efb9e19094440c317e21b-Abstract-round1.html}{CUAD: An Expert-Annotated NLP Dataset for Legal Contract Review}. In Advances in Neural Information Processing Systems (NeurIPS), 2021.

\bibitem{koreeda2021contractnli}
\textit{Y. Koreeda and C. D. Manning},
\href{https://aclanthology.org/2021.findings-emnlp.164/}{ContractNLI: A Dataset for Document-level Natural Language Inference for Contracts}. Findings of the Association for Computational Linguistics: EMNLP, 2021.

\bibitem{dominguez2024lawma}
\textit{R. Dominguez-Olmedo, V. Nanda, R. Abebe, S. Bechtold, C. Engel, et al.},
\href{https://arxiv.org/abs/2407.16615}{Lawma: The Power of Specialization for Legal Tasks}. arXiv preprint arXiv:2407.16615, 2024.

\bibitem{colombo2024saullm}
\textit{P. Colombo, T. P. Pires, M. Boudiaf, D. Culver, R. Melo, C. Corro, A. F. T. Martins, F. Esposito, V. L. Raposo, S. Morgado, and M. Desa},
\href{https://arxiv.org/abs/2403.03883}{SaulLM-7B: A pioneering Large Language Model for Law}. arXiv preprint arXiv:2403.03883, 2024.

\bibitem{cui2023chatlaw}
\textit{J. Cui, Z. Li, Y. Yan, B. Chen, and L. Yuan},
\href{https://www.researchgate.net/publication/371943604_ChatLaw_Open-Source_Legal_Large_Language_Model_with_Integrated_External_Knowledge_Bases}{ChatLaw: Open-Source Legal Large Language Model with Integrated External Knowledge Bases}. arXiv preprint, 2023.

\bibitem{kalra2024hypa}
\textit{R. Kalra, Z. Wu, A. Gulley, A. Hilliard, X. Guan, A. Koshiyama, and P. Treleaven},
\href{https://doi.org/10.48550/arxiv.2409.09046}{HyPA-RAG: A Hybrid Parameter Adaptive Retrieval-Augmented Generation System for AI Legal and Policy Applications}. arXiv:2409.09046, 2024.

\bibitem{barron2025bridging}
\textit{R. C. Barron, M. E. Eren, O. M. Serafimova, C. Matuszek, and B. S. Alexandrov},
\href{https://doi.org/10.48550/arxiv.2502.20364}{Bridging Legal Knowledge and AI: Retrieval-Augmented Generation with Vector Stores, Knowledge Graphs, and Hierarchical Non-negative Matrix Factorization}. arXiv:2502.20364, 2025.

\bibitem{pipitone2024legalbench}
\textit{N. Pipitone and G. H. Alami},
\href{https://doi.org/10.48550/arxiv.2408.10343}{LegalBench-RAG: A Benchmark for Retrieval-Augmented Generation in the Legal Domain}. arXiv:2408.10343, 2024.

\bibitem{li2025lexrag}
\textit{H. Li, Y. Chen, Y. Hu, Q. Ai, J. Chen, X. Yang, et al.},
\href{https://doi.org/10.48550/arxiv.2502.20640}{LexRAG: Benchmarking Retrieval-Augmented Generation in Multi-Turn Legal Consultation Conversation}. arXiv:2502.20640, 2025.

\bibitem{thanh2024krag}
\textit{N. H. Thanh and K. Satoh},
\href{https://doi.org/10.48550/arxiv.2410.07551}{KRAG Framework for Enhancing LLMs in the Legal Domain}. arXiv:2410.07551, 2024.

\bibitem{xi2023rise}
\textit{Z. Xi, W. Chen, X. Guo, W. He, Y. Ding, B. Hong, et al.},
\href{https://arxiv.org/abs/2309.07864}{The Rise and Potential of Large Language Model Based Agents: A Survey}. arXiv preprint arXiv:2309.07864, 2023.

\bibitem{wang2024agent}
\textit{J. Wang and Z. Duan},
\href{https://doi.org/10.48550/arxiv.2412.03801}{Agent AI with LangGraph: A Modular Framework for Enhancing Machine Translation Using Large Language Models}. arXiv:2412.03801, 2024.

\bibitem{rauba2026multi}
\textit{P. Rauba, S. Cepenas, and M. van der Schaar},
\href{https://arxiv.org/abs/2601.23211}{Multi-Agent Systems Should be Treated as Principal-Agent Problems}. arXiv preprint arXiv:2601.23211, 2026.

\bibitem{chalkidis2022lexglue}
\textit{I. Chalkidis, A. Jana, D. Hartung, M. Bommarito, I. Androutsopoulos, D. M. Katz, and N. Aletras},
\href{https://aclanthology.org/2022.acl-long.297/}{LexGLUE: A Benchmark Dataset for Legal Language Understanding in English}. In Proceedings of the 60th Annual Meeting of the Association for Computational Linguistics (ACL), 2022.

\bibitem{guha2023legalbench}
\textit{N. Guha, J. Nyarko, D. Ho, C. R{\'e}, A. Chilton, et al.},
\href{https://papers.neurips.cc/paper_files/paper/2023/hash/89e44582fd28ddfea1ea4dcb0ebbf4b0-Abstract-Datasets_and_Benchmarks.html}{LegalBench: A Collaboratively Built Benchmark for Measuring Legal Reasoning in Large Language Models}. In Advances in Neural Information Processing Systems (NeurIPS), 2023.

\end{thebibliography}

\end{document}
